% ============================================================================
%  SECTION 6 — Equilibrated constrained mixture
% ============================================================================
\begin{frame}[t]{}
  \vfill\begin{center}{\Huge\color{YaleBlue} Equilibrated CMM}
  \end{center}\vfill
\end{frame}

\begin{frame}[t]{Skip the Transient --- Solve for the End State}
  Set every rate to zero and solve the \textbf{algebraic} equilibrium --- one
  root-find instead of thousands of time steps.
  \dlab{general form --- the 3D equilibrium conditions}
  \vspace{-0.2em}
  \giveneq{\underbrace{\Upsilon=1\Rightarrow\bar\sigma^*=\bar\sigma_h}_{\text{(A) growth balance}}
    \quad
    \underbrace{\bm{F}_n^{k}=\text{const}\Rightarrow\sigma^k=\sigma_h^k}_{\text{(B) remodeling done}}
    \quad
    \underbrace{\nabla\!\cdot\bm{\sigma}=\bm 0}_{\text{(C) mechanical equilib.}}}
  \zlab{reduction --- scalar conditions at the evolved geometry}
  \giveneq{\bar\sigma^*=\bar\sigma_h,\qquad \lambda_e^k=G^k,\qquad
    \sigma_\theta=\tfrac{P a}{h}.}
  \bigfigw{equil_longterm.png}{0.56\linewidth}{The equilibrated model jumps
  straight to the long-term evolved state that the transient theories approach.}
  \srccite{Latorre, Humphrey, ZAMM (2018)}
\end{frame}

\begin{frame}[t]{Collapse to One Scalar Equation}
  Let the turnover constituents grow in fixed proportion (exact when they share
  gain and turnover). The 3D conditions (A)--(C) collapse to a \emph{single}
  equation for the evolved stretch $\lambda^*$ --- consistent with the same
  $\lambda,\ G^k,\ \phi^k$ used throughout:
  \zlab{scalar equilibrium equation for $\lambda^*$}
  \slboxeq{2.8em}{eq:equil}{\phi^e\,s
    + s\,\frac{\bar\sigma_h - \sigma^e(G^e\lambda^*)}{\sigma_{\text{turn}} - \sigma_h^e}
    - \gamma\,(\lambda^*)^{n} = 0}
  with $n$ the geometry exponent ($n=1$ bar, $n=2$ artery), $\gamma$ the sustained
  load factor, $s$ the surviving elastin fraction, $\phi^e$ the elastin mass
  fraction. Once $\lambda^*$ is known, the evolved mass ratio is
  $M/M_0=\gamma(\lambda^*)^n$.
  \vspace{0.3em}

  This equation \emph{is} the fixed point of the homogenized ODEs
  (\eqref{eq:homog_mass}--\eqref{eq:homog_remodel} with $\dd/\dd t=0$), so its root
  \textbf{coincides with the long-time limit of the transient theories} --- to
  several digits.
  \srccite{Latorre, Humphrey, ZAMM (2018)}
\end{frame}

\begin{frame}[t]{The Punchline: Existence = Stability}
  Equation \eqref{eq:equil} may have \textbf{no physical root} --- then no adapted
  state exists and the transient theories dilate without bound. The solver is
  therefore the cleanest \textbf{stability test}: root $\Rightarrow$ adapts, no root
  $\Rightarrow$ unstable. Near the boundary the transient slows critically, which
  is exactly why the instant algebraic test is so useful.
\end{frame}


% ---- interactive coding exercise: ex05 -------------------------------------
\begin{frame}[t]{}
  \begin{yourturnbox}{Equilibrated --- code}
    \qq{Check that the algebraic equilibrium matches the transient end-state --- then
    find the elastin loss at which the equilibrium \emph{disappears}.
    \qkind{\texttt{ex05}}}
    \qres{\texttt{exercises/ex05\_equilibrated.py}. Tunable:
    \texttt{ELASTIN\_SURVIVING=0.3} (try 0.4 / 0.3 / 0.2).}
    \qhint{Part A: the equilibrated $\lambda^*$ should sit on the transient
    plateau. Part B: lower $s$ until the root vanishes --- that is $s_{\text{crit}}$;
    run just below it and watch the transient diverge.}
    \qrun{uv run python exercises/ex05\_equilibrated.py}
  \end{yourturnbox}
\end{frame}

\begin{frame}[t]{}
  \begin{answerbox}{Equilibrated}
    Where a root exists, the equilibrated solve lands exactly on the transient
    end-state. Below $s_{\text{crit}}$ there is \emph{no} root --- the transient
    dilates without bound.
  \end{answerbox}
  \vspace{-0.3em}
  \bigfigw{fig05_equilibrated.pdf}{0.66\linewidth}{Your run: equilibrated end-state
  on the transient plateau where a root exists; near the boundary, critical slowing.}
  \vspace{-0.4em}
  \begin{center}{\footnotesize\color{ABlue}\bfseries Learning:} \textbf{matches} the
  transient theories \emph{if} an equilibrium exists --- and flags when it does not.\end{center}
\end{frame}
